% =====================================================================
%  Chapter 26 — Kelly with an Uncertain Mean
% =====================================================================

\chapter{Kelly with an Uncertain Mean}
\label{ch:kelly}

The book's sizing rule is a Kelly criterion applied to an uncertain
forward return. This chapter derives the criterion, shows why the standard
form cannot be used with an estimated mean, and derives the
certainty-equivalent form the book actually uses.

\section{The criterion}

Suppose a strategy wins a fraction \(p\) of its bets and loses the rest,
with the win paying \(b\) times the stake. The growth rate of the
investor's wealth, per bet, is
\begin{equation}
g(p) = p\log(1+pb) - (1-p)\log(1-p),
\end{equation}
and the maximisation of \(g\) gives
\begin{equation}
p^{*} = \frac{pb - (1-p)}{b},
\qquad\text{equivalently}\qquad
f^{*} = \frac{pb - (1-p)}{b},
\end{equation}
where \(f^{*}\) is the fraction of wealth to stake. The well-known
property of this rule is that it maximises the expected \emph{log}
wealth, which means it asymptotically beats every other rule --- provided
the probabilities are correct.

The continuous-time analogue, for a risky asset with excess expected return
\(\mu\) and volatility \(\sigma\), is the continuous-time Kelly rule
applied to a diffusion, which gives the familiar
\begin{equation}
f^{*} = \frac{\mu}{\sigma^{2}},
\end{equation}
the fraction of wealth in the risky asset under logarithmic utility with a
quadratic penalty. And for a portfolio of several assets with excess means
\(\mu\) and covariance \(\Sigma\), the mean-variance Kelly solution is
\begin{equation}
w^{*} = \Sigma^{-1}\mu.
\end{equation}

The attraction of the criterion is that it is the only rule that
compounds. A portfolio that runs at half Kelly gives up a little growth in
exchange for a materially lower probability of ruin, which is the trade
Chapter 7 said a saver should care about.

\begin{takeaway}[The attraction and the trap]
Kelly maximises expected log wealth, so it beats every other rule in the
long run --- if the probabilities are right. The rest of this chapter is
about what happens when they are not.
\end{takeaway}

\section{Why the standard form cannot be used}

The rule above assumes \(\mu\) is known. In this book it is not: it is a
posterior mean with a posterior variance attached, and Chapter 20 showed
that the posterior variance is large relative to the posterior mean.

This is not a technical inconvenience; it is the central fact of the
problem. When \(\mu\) is estimated with error, the Kelly rule computed
from \(\hat{\mu}\) is not merely noisy but \emph{biased upward in its
variance}, and a strategy that runs full Kelly on an uncertain mean can
have a negative expected log growth rate --- the well-known result that
overbetting on a noisy signal loses money at an exponential rate.

The remedy is standard and appears in the literature in several forms. The
common thread is to penalise the estimate by its own uncertainty. The
book uses the certainty-equivalent construction.

\section{The certainty-equivalent construction}

Let the posterior over the excess mean of asset \(i\) be
\(N(\mu_i, \psi_i)\), with \(\psi_i\) the posterior variance of the mean.
Under a quadratic utility with relative risk aversion \(\gamma\), the
certainty-equivalent excess return is
\begin{equation}
\tilde{\mu}_i = \mu_i - \frac{\gamma}{2}\psi_i,
\end{equation}
because the mean of a concave function of a random variable is less than
the function of the mean, and to first order the correction is exactly the
penalty term above.

Substituting into the mean-variance portfolio problem gives the book its
sizing rule:
\begin{equation}
w^{CE} = \frac{1}{\gamma}\left(\Sigma + \Psi\right)^{-1}\mu,
\qquad
\Psi = \operatorname{diag}(\psi_1,\dots,\psi_N),
\end{equation}
where the \(\Psi\) term is added to the covariance because the
uncertainty in the mean is a risk that the investor bears when sizing.

The formula has four consequences, and each of them is the reason it is
used.

\textbf{The position falls as the doubt rises.} If \(\psi_i\) grows, the
diagonal of \(\Sigma+\Psi\) grows, and the position in asset \(i\) falls.
A forecast the model is unsure about produces a small trade.

\textbf{The position is bounded in aggregate.} The sum of the
certainty-equivalent weights is smaller than the sum of the raw ones, so
the tilt cannot be larger than the raw signal implies.

\textbf{There is a natural closed form for the shrinkage.} For a single
asset with a signal-to-noise ratio \(z = \mu/\sigma\), the position falls
with \(z\) approximately as
\begin{equation}
w \propto \frac{z}{1 + z^{2}},
\end{equation}
which is maximal at \(z = 1\) --- when the forecast is one standard
deviation from zero --- and then falls again. This is the classical
result that a signal must be both positive and strong to be worth
trading, and it is the single most important shape in the book.

\textbf{The scaling is explicit.} The rule is linear in the forecast and
quadratic in the doubt, and the book's parameter --- the fraction of full
Kelly --- multiplies the whole thing.

How the position size responds to the forecast and to the uncertainty in it is shown in \figref{fig:kelly-shrinkage}.

\section{The fraction of Kelly}

The book's default is a quarter of full Kelly, and the reason is not
caution for its own sake. It is that the \emph{estimated} mean is a noisy
estimate of the true mean, and the certainty-equivalent correction
already accounts for the posterior variance --- but the posterior variance
is itself an estimate, and the additional shrinkage of a fractional Kelly
is cheap insurance against that estimate being wrong.

The two most common arguments against fractional Kelly are worth
addressing. The first is that it gives up a quarter of the long-run edge.
The answer is that the edge being given up is an edge that is being
estimated, and a strategy that runs full Kelly on an estimate has a
materially higher probability of ruin --- the probability that the
book's own ruin calculation in Chapter 7 is there to make visible. The
second is that the optimal fraction depends on the accuracy of the
estimate, so any fixed fraction is arbitrary. The answer is that it does,
and that the book's trial family includes fractions of 0.10, 0.25 and
0.50 so that the reader can see what the choice is worth.

\begin{keytakeaways}
\item Kelly maximises expected log wealth and beats every other rule in
the long run, provided the probabilities are right.
\item With an estimated mean, the standard rule overbets, and the
certainty-equivalent construction penalises the estimate by its own
posterior variance: \(w = \gamma^{-1}(\Sigma+\Psi)^{-1}\mu\).
\item The position falls as the doubt rises and is maximised at a
signal-to-noise ratio near one, so a weak signal produces no trade.
\item The book's default is a quarter of full Kelly, and the fraction is
one of the trial-family parameters.
\end{keytakeaways}
